\documentclass[11pt,a4paper]{report}

% ═══════════════════════════════════════════════════════════════
%  Dossier: Vacancy tax / anti-grandes poseedores
%  Public Economics & Political Economy — BIEBAM Segovia
% ═══════════════════════════════════════════════════════════════

\usepackage[T1]{fontenc}
\usepackage[utf8]{inputenc}
\usepackage[spanish,es-tabla]{babel}
\usepackage{lmodern}
\usepackage[margin=2.2cm,top=2.6cm,bottom=2.6cm,headheight=22pt]{geometry}
\usepackage{microtype}
\usepackage{amsmath,amssymb}
\usepackage{booktabs}
\usepackage{longtable}
\usepackage{array}
\usepackage{tabularx}
\usepackage{enumitem}
\usepackage{xcolor}
\usepackage{graphicx}
\usepackage{tikz}
\usepackage{fancyhdr}
\usepackage{titlesec}
\usepackage{hyperref}
\usepackage{csquotes}
\usepackage[most]{tcolorbox}
\usepackage{parskip}
\usepackage{caption}
\usepackage{multirow}
\usepackage{float}

\hypersetup{
  colorlinks=true,
  linkcolor=azuloscuro,
  urlcolor=azulmedio,
  citecolor=azuloscuro,
  pdftitle={Dossier Vacancy Tax --- Public Economics Housing},
  pdfauthor={Equipo Housing / Vacancy Tax}
}

% ── Colores ──────────────────────────────────────────────────
\definecolor{azuloscuro}{HTML}{1B3A4B}
\definecolor{azulmedio}{HTML}{2E6B8A}
\definecolor{verdeint}{HTML}{1F6B4A}
\definecolor{ambar}{HTML}{B45309}
\definecolor{rojoaviso}{HTML}{9B1C1C}
\definecolor{fondosuave}{HTML}{F4F7FA}
\definecolor{fondoverde}{HTML}{ECF6F0}
\definecolor{fondoambar}{HTML}{FFF7ED}
\definecolor{fondorojo}{HTML}{FEF2F2}
\definecolor{linea}{HTML}{CBD5E1}

% ── Boxes ────────────────────────────────────────────────────
\newtcolorbox{intuicion}[1][]{%
  enhanced, breakable,
  colback=fondoverde, colframe=verdeint,
  boxrule=0.6pt, arc=2pt,
  left=10pt, right=10pt, top=8pt, bottom=8pt,
  fonttitle=\bfseries\color{verdeint},
  title={Intuici\'on},
  #1
}

\newtcolorbox{modelo}[1][]{%
  enhanced, breakable,
  colback=fondosuave, colframe=azulmedio,
  boxrule=0.6pt, arc=2pt,
  left=10pt, right=10pt, top=8pt, bottom=8pt,
  fonttitle=\bfseries\color{azuloscuro},
  title={Modelo (versi\'on sencilla)},
  #1
}

\newtcolorbox{dato}[1][]{%
  enhanced, breakable,
  colback=fondoambar, colframe=ambar,
  boxrule=0.6pt, arc=2pt,
  left=10pt, right=10pt, top=8pt, bottom=8pt,
  fonttitle=\bfseries\color{ambar},
  title={Dato ancla},
  #1
}

\newtcolorbox{cuidado}[1][]{%
  enhanced, breakable,
  colback=fondorojo, colframe=rojoaviso,
  boxrule=0.6pt, arc=2pt,
  left=10pt, right=10pt, top=8pt, bottom=8pt,
  fonttitle=\bfseries\color{rojoaviso},
  title={Cuidado / no exagerar},
  #1
}

\newtcolorbox{frase}[1][]{%
  enhanced, breakable,
  colback=white, colframe=azuloscuro,
  boxrule=1pt, arc=0pt,
  left=12pt, right=12pt, top=10pt, bottom=10pt,
  fontupper=\itshape\large,
  #1
}

\newtcolorbox{resumenbox}[1][]{%
  enhanced, breakable,
  colback=azuloscuro!5, colframe=azuloscuro,
  boxrule=0.8pt, arc=2pt,
  left=10pt, right=10pt, top=8pt, bottom=8pt,
  fonttitle=\bfseries\color{azuloscuro},
  title={En una frase},
  #1
}

% ── Headers ───────────────────────────────────────────────────
\pagestyle{fancy}
\fancyhf{}
\fancyhead[L]{\small\color{azulmedio} Vacancy tax --- dossier de trabajo}
\fancyhead[R]{\small\color{azulmedio} Public Economics}
\fancyfoot[C]{\thepage}
\renewcommand{\headrulewidth}{0.4pt}
\renewcommand{\headrule}{\hbox to\headwidth{\color{linea}\leaders\hrule height \headrulewidth\hfill}}

\titleformat{\chapter}[display]
  {\normalfont\huge\bfseries\color{azuloscuro}}
  {\chaptertitlename\ \thechapter}{12pt}{\Huge}
\titleformat{\section}
  {\normalfont\Large\bfseries\color{azuloscuro}}
  {\thesection}{1em}{}
\titleformat{\subsection}
  {\normalfont\large\bfseries\color{azulmedio}}
  {\thesubsection}{1em}{}

\setcounter{tocdepth}{2}
\setlist[itemize]{leftmargin=1.4em, itemsep=0.25em}
\setlist[enumerate]{leftmargin=1.6em, itemsep=0.25em}

\begin{document}

% ══════════════════════════════════════════════════════════════
\begin{titlepage}
\raggedright
\vspace*{1.5cm}
{\color{azulmedio}\rule{\textwidth}{2pt}}\\[0.6cm]
{\fontsize{28}{34}\selectfont\bfseries\color{azuloscuro}
  Impuesto a viviendas vac\'ias\\[0.2cm]
  Dossier de argumentos, modelos y evidencia}\\[0.5cm]
{\color{azulmedio}\rule{\textwidth}{1.2pt}}\\[1.2cm]
{\Large\color{azuloscuro}
  Public Economics \& Political Economy\\[0.3cm]
  Debate \emph{Government vs Parliament} --- Tema Housing}\\[0.8cm]
{\large
  Posici\'on: \textbf{gobierno fiscal / anti-grandes poseedores}\\[0.2cm]
  Pol\'itica estrella: \textbf{vacancy tax} sobre vacancy prolongada\\[0.2cm]
  Fecha del debate: 1 de diciembre de 2026}\\[1.5cm]
\begin{tcolorbox}[colback=fondosuave, colframe=azulmedio, boxrule=0.5pt, arc=2pt]
  Documento de trabajo para el equipo. Recoge \textbf{todo} lo elaborado hasta ahora:
  posici\'on, argumentos del curso, modelos de los papers, evidencia causal,
  datos oficiales (BdE, INE, Catastro, MIVAU), dise\~no de la propuesta,
  respuestas al equipo construction-led, y trampas a evitar.
  Estilo: intuitivo primero; f\'ormulas solo cuando ayudan.
\end{tcolorbox}
\vfill
{\small\color{azulmedio}
  Fuentes locales en \texttt{docs/}, \texttt{references/}, \texttt{data/}, \texttt{research/briefs/}.\\
  Compilado a partir del material del repositorio PublicEconomicsHousing.}
\end{titlepage}

\tableofcontents
\newpage

% ══════════════════════════════════════════════════════════════
\chapter{Qui\'enes somos y qu\'e defendemos}
\label{ch:posicion}

\section{La pregunta del debate}

El syllabus plantea dos v\'ias frente a la crisis de vivienda en Espa\~na:

\begin{quote}
\emph{Should we tackle housing affordability primarily through expanding construction
or through taxation aimed at discouraging speculative ownership?}
\end{quote}

No son mutuamente excluyentes. Pero obligan a elegir una \textbf{prioridad}.
Nosotros somos el gobierno de la v\'ia fiscal.

\begin{center}
\renewcommand{\arraystretch}{1.25}
\begin{tabular}{@{}ll@{}}
\toprule
\textbf{Campo} & \textbf{Decisi\'on} \\
\midrule
Tema & 1 --- \emph{The Housing Crisis: Who Should Build and Own?} \\
Lado & Tributaci\'on (no expansi\'on de oferta v\'ia construcci\'on) \\
Encuadre & Anti-grandes poseedores / anti-especulaci\'on \\
Pol\'itica & Impuesto a viviendas vac\'ias (\emph{vacancy tax}) \\
\bottomrule
\end{tabular}
\end{center}

\section{Narrativa en 30 segundos}

\begin{frase}
La asequibilidad en Espa\~na no se resuelve solo construyendo m\'as.
Hay viviendas que ya existen y no se usan como hogar.
Nuestro trabajo es corregir los incentivos que permiten mantenerlas vac\'ias
como activo especulativo.
\end{frase}

\section{Formato del d\'ia}

\begin{itemize}
  \item Presentaci\'on del gobierno: \textbf{10 minutos}.
  \item Preguntas del parlamento: \textbf{6 minutos}.
  \item Voto + cambio de grupo: \textbf{2 minutos}.
  \item Todos hablan; claims con teor\'ia + evidencia; presentaci\'on en ingl\'es.
\end{itemize}

\section{Qu\'e pide el brief (checklist)}

Antes de salir a escena, la propuesta debe dejar claro:

\begin{enumerate}
  \item Definici\'on de \textbf{vivienda vac\'ia} (tiempo, registros, excepciones).
  \item \textbf{Qui\'en paga} (\'enfasis en grandes poseedores / umbrales).
  \item Base, tipo(s), exenciones (obras, herencia reciente,\ldots).
  \item \textbf{Enforcement} (catastro, utilities, cruces administrativos).
  \item Uso de la recaudaci\'on (vivienda p\'ublica / limited-profit).
  \item Canal causal: vac\'ia $\downarrow$ $\Rightarrow$ oferta efectiva $\uparrow$ $\Rightarrow$ acceso.
  \item Incidencia y winners/losers.
  \item Respuestas al contraargumento construction-led.
  \item Ancla emp\'irica: \textbf{Seg\'u (2020)} + datos ES (BdE / INE).
\end{enumerate}

% ══════════════════════════════════════════════════════════════
\chapter{La historia completa en una cadena}
\label{ch:cadena}

Antes de entrar argumento por argumento, conviene tener el hilo entero
en la cabeza. Es la cadena causal que queremos contar en 10 minutos.

\begin{center}
{\shorthandoff{<>}%
\begin{tikzpicture}[
  box/.style={draw=azulmedio, fill=fondosuave, rounded corners=2pt,
              align=center, text width=11.5cm, inner sep=7pt, font=\small},
  flecha/.style={-latex, thick, color=azulmedio}
]
\node[box] (a) at (0,0) {\textbf{Demanda alta} + oferta nueva lenta (BdE)\\
  + stock desviado: vac\'ios, 2\textsuperscript{a}s, tur\'istico (BdE + Airbnb BCN)\\
  + incentivos privados a retener vac\'io (opci\'on real / holding estrat\'egico)};
\node[box] (b) at (0,-2.1) {\textbf{Fallo:} inactive housing + externalidades + equity};
\node[box] (c) at (0,-3.8) {\textbf{Vacancy tax} eleva el coste de no participar};
\node[box] (d) at (0,-5.5) {Vacancy inactiva $\downarrow$ $\rightarrow$ m\'as residencias principales\\
  (evidencia Francia: $-13\%$; Seg\'u 2020)};
\node[box] (e) at (0,-7.2) {Oferta efectiva \textbf{r\'apida} + recaudaci\'on $\rightarrow$ vivienda asequible\\
  Incidencia sobre carteras grandes inactivas};
\draw[flecha] (a) -- (b);
\draw[flecha] (b) -- (c);
\draw[flecha] (c) -- (d);
\draw[flecha] (d) -- (e);
\end{tikzpicture}}
\end{center}

\begin{resumenbox}
Misma crisis que el otro gobierno; distinto cuello de botella a corto plazo:
no solo ``faltan permisos'', sino \textbf{incentivos de uso del stock}.
La construcci\'on es complemento; la vacancy tax es nuestra prioridad.
\end{resumenbox}

% ══════════════════════════════════════════════════════════════
\chapter{Argumentos del curso (hablar el idioma del profesor)}
\label{ch:curso}

El PDF del curso (Weeks 1--10, 12 + PS1--8) no tiene un paper de vacancy tax.
Lo que s\'i tiene es el \textbf{andamiaje te\'orico} para justificar la medida
y anticipar objeciones. Aqu\'i cada pieza va con intuici\'on primero.

\section{N\'ucleo: eficiencia + equity}

\subsection{A1. Primer teorema de bienestar + derechos incompletos}

\begin{intuicion}
El mercado ``funciona bien'' solo si nadie impone costes a terceros que no
paga. Si dejar un piso vac\'io en un barrio tensionado perjudica a inquilinos
y vecinos, el equilibrio laissez-faire \textbf{no} es eficiente.
\end{intuicion}

\textbf{Qu\'e dice el curso.} El First Fundamental Theorem exige, entre otras cosas,
\emph{no externalities / perfectly defined property rights}. Si falla, el
equilibrio no es Pareto eficiente. Closing line de externalidades (Week~4):
\emph{``Externalities arise because of incomplete property rights.''}

\textbf{Extrapolaci\'on.} El t\'itulo de propiedad puede estar claro; el
\textbf{uso} (vac\'io prolongado) genera efectos no cotizados: menos oferta
efectiva, presi\'on en rentas, degradaci\'on urbana.

\begin{frase}
``El 1.\textsuperscript{er} teorema no aplica: hay un gap entre el coste privado
de dejar vac\'io y el coste social de restar una vivienda al mercado residencial.''
\end{frase}

\subsection{A2. Pigou: impuesto correctivo = da\~no marginal}

\begin{intuicion}
Si contaminar es malo, no hace falta prohibir: se pone un precio al da\~no.
Aqu\'i el ``contaminante'' es mantener vac\'io prolongado. El impuesto no es
castigo gen\'erico a la propiedad; es un precio sobre la actividad da\~nina.
\end{intuicion}

\begin{modelo}
En el \'optimo, el tipo Pigouviano cumple $t^{*} =$ da\~no marginal social.
Analog\'ia limpia con el carbon tax del PS4: el laissez-faire produce
``demasiado'' del mal (emisiones / vac\'ios).
\end{modelo}

\textbf{Dise\~no Weitzman / PS4.} Sin incertidumbre fuerte sobre el da\~no
marginal, el instrumento de \textbf{precio} es natural frente a cupos r\'igidos
de ocupaci\'on forzada. Si el parlamento pide ``expropiar'', Weitzman os deja
defender el tax.

\textbf{L\'imite Coase.} El bargaining privado no escala cuando hay muchos
afectados (igual que contaminaci\'on atmosf\'erica). Los inquilinos potenciales
no pueden negociar con cada gran poseedor de vac\'ios.

\subsection{A3. Equity no es efficiency; Sen; segundo teorema}

\begin{intuicion}
Un mercado puede ser ``eficiente'' y a la vez socialmente inaceptable.
Pareto no es suficiente. Adem\'as, en la pr\'actica no podemos redistribuir
vivienda con lump-sum m\'agicos: hay trade-off equity--efficiency.
\end{intuicion}

\textbf{Qu\'e dice el curso.}
\begin{itemize}
  \item \emph{``A society can be Pareto optimal and still perfectly disgusting''} (Sen).
  \item SFTWE: cualquier Pareto se alcanza redistribuyendo endowments y dejando al mercado;
        lump-sum no es factible $\Rightarrow$ trade-off.
  \item Week~12: intervenci\'on por (1)~fallos de mercado y (2)~redistribuci\'on.
\end{itemize}

\textbf{Uso.} Incluso si el oponente dice ``el mercado constructivo es eficiente'',
respond\'eis en equity: concentraci\'on de stock + vac\'ios puede ser
Pareto-compatible y socialmente inaceptable. La vacancy tax ataca uso del stock
\textbf{y} carga sobre multi-propietarios.

\subsection{A4. Capital / wealth m\'as concentrados que labour}

\begin{intuicion}
Los ingresos del trabajo est\'an menos concentrados que la riqueza.
La vivienda es un veh\'iculo t\'ipico de esa riqueza. El problema no es solo
``faltan pisos nuevos'': es \textbf{qui\'en controla el stock} y con qu\'e incentivos.
\end{intuicion}

\textbf{Hechos de clase.} Top~10\% labour $\sim$20--30\% vs capital $\sim$50--90\%;
wealth m\'as concentrado que income.

\section{Dise\~no, incidencia y medici\'on}

\subsection{B5. Incidencia econ\'omica frente a statutory incidence}

\begin{intuicion}
Quien ``aparece'' en la ley como pagador no es necesariamente quien pierde
poder adquisitivo. Depende de elasticidades. Anticipad: ``el impuesto lo pagan
los inquilinos''.
\end{intuicion}

\textbf{Respuesta de clase.} Con stock r\'igido y multi-propietarios que pueden
poner en alquiler, parte de la carga cae en owners. El canal deseado es
\textbf{cambio de comportamiento} (dejar de vaciar), no recaudar por recaudar.
El curso usa Linden \& Rockoff (2008) como ejemplo expl\'icito de housing.

\subsection{B6. Capitalizaci\'on en precios}

\begin{intuicion}
Las mejoras urbanas se ``meten'' en el precio de la vivienda. Si no grav\'ais
ownership especulativo, los gains de pol\'itica p\'ublica los capturan
landowners --- otra raz\'on para gravar usos vac\'ios/especulativos.
\end{intuicion}

El curso usa house prices como m\'etrica de bienestar: Linden--Rockoff,
Chay--Greenstone, Samuelson (revelaci\'on de preferencias), CBA Week~7.

\subsection{B7. Elasticidades y DWL / ``ETI inmobiliaria''}

\begin{intuicion}
Un impuesto falla si la gente responde cambiando de etiqueta (evasión) en
lugar de cambiando de comportamiento real. Por eso importan definici\'on,
enforcement y umbrales.
\end{intuicion}

\textbf{Checklist de dise\~no.} Definici\'on operativa de ``vac\'ia'', cruces
catastro/utilities, umbrales para grandes poseedores $\rightarrow$ bajar la
evasión (Airbnb, sociedades pantalla,\ldots). Week~12: si la base es el\'astica
por redefinici\'on, ensanchad la base y reforzad enforcement.

\subsection{B8. Timing: stock vs flujo (PS7)}

\begin{intuicion}
Construir da beneficios tarde (permisos, obras). Poner vac\'ios en uso da
beneficios antes. Con descuento alto o impaciencia pol\'itica, el instrumento
sobre stock gana el CBA de timing --- exactamente la l\'ogica del PS7 con dos
housing projects.
\end{intuicion}

\section{Econom\'ia pol\'itica (parlamento)}

\subsection{C9. Costes concentrados, beneficios difusos}

\begin{intuicion}
Una vacancy tax seria tiene lobby de grandes poseedores; los beneficiarios
(inquilinos) est\'an dispersos. El rechazo pol\'itico no prueba ineficiencia
econ\'omica (Week~12, Prop~40 / wealth tax California).
\end{intuicion}

\subsection{C10. Tax-and-transfer + median voter}

\begin{intuicion}
Si el votante mediano es ``pobre'' relativo a la media de endowments, puede
preferir gravar ownership especulativo. No afirm\'eis que el MVT ``prueba''
la vacancy tax; s\'i que la desigualdad inmobiliaria entra en el equilibrio pol\'itico.
\end{intuicion}

\subsection{C11. Political Coase Theorem falla}

\begin{intuicion}
No esper\'eis un acuerdo voluntario limpio entre grandes poseedores y
households sin techo asequible. Hace falta \textbf{regla fiscal} con enforcement,
no solo promesas de ``poner en alquiler'' (Acemoglu 2003 / Week~6).
\end{intuicion}

\subsection{C12. ``Los ricos se van'' --- magnitud vs an\'ecdota}

\begin{intuicion}
El curso (Jakobsen et al.\ sobre wealth tax) ense\~na a separar respuesta
migratoria de owners vs efecto macro. Pedid evidencia de magnitud.
Una vacancy tax acotada a vacancy prolongada es menos ``exit-inducing''
que un wealth tax amplio.
\end{intuicion}

\section{Analog\'ias \'utiles (con matices)}

\begin{center}
\renewcommand{\arraystretch}{1.2}
\begin{tabularx}{\textwidth}{@{}>{\raggedright\arraybackslash}p{3.2cm}
  >{\raggedright\arraybackslash}X
  >{\raggedright\arraybackslash}X@{}}
\toprule
\textbf{Tema de clase} & \textbf{Analog\'ia cuidadosa} & \textbf{No dig\'ais} \\
\midrule
Public goods / free riding &
Asequibilidad urbana tiene rasgos de bien p\'ublico local &
Que la vivienda sea un pure public good \\
Lindahl / preference revelation &
Capitalizaci\'on y datos de vac\'ios ayudan a calibrar &
Que conocemos el MD con precisi\'on quir\'urgica \\
Moral hazard &
Subsidiar solo oferta nueva sin gravar vac\'io puede ``asegurar'' el hold-out &
Que vacancy = moral hazard de UI \\
Crowd-out &
Inversi\'on p\'ublica puede desplazarse si el stock privado se retira a vac\'io/tur\'istico &
Equivalencia exacta con carbon offsets \\
Piketty $r$ vs $g$ &
Capital (incl.\ housing) captura renta; el tax reduce el return a \emph{no usar} &
Que $r>g$ ``demuestra'' vacancy tax \\
\bottomrule
\end{tabularx}
\end{center}

\section{Guion m\'inimo de teor\'ia (unos 2 minutos)}

\begin{enumerate}
  \item \textbf{Diagn\'ostico:} FFTWE falla --- coste social $\neq$ coste privado.
  \item \textbf{Equity:} capital/housing wealth concentrados; Pareto puede ser ``disgusting''.
  \item \textbf{Instrumento:} Pigou sobre vacancy prolongada; Coase privado no escala.
  \item \textbf{Dise\~no:} base, enforcement, bajar evasión; ingresos a vivienda asequible.
  \item \textbf{Vs construction-only:} timing --- stock responde antes que obra nueva.
  \item \textbf{Incidencia / capitalizaci\'on:} owners vs renters; qui\'en captura land rents.
  \item \textbf{Pol\'itica:} costes concentrados; migraci\'on $\neq$ colapso macro.
\end{enumerate}

Luego ancl\'ais evidencia: Seg\'u (2020) + BdE + Airbnb Barcelona si habl\'ais de desv\'io de uso.

% ══════════════════════════════════════════════════════════════
\chapter{N\'ucleo te\'orico desde los papers: por qu\'e hay vac\'ios}
\label{ch:teoria-papers}

\section{Dos tipos de vacancy (la idea m\'as importante)}

\begin{intuicion}
No todos los vac\'ios son ``especulaci\'on''. Hay vac\'ios normales de b\'usqueda
(como el paro friccional) y vac\'ios voluntarios: el due\~no \textbf{elige}
no participar. Nuestro impuesto apunta al segundo, no al primero.
\end{intuicion}

\begin{modelo}[title={Modelo (versi\'on sencilla) --- Seg\'u / Wheaton}]
Dos componentes:
\begin{enumerate}
  \item \textbf{Frictional / open vacancies ($v$).}
        El propietario \emph{quiere} alquilar o vender, pero tarda en emparejar
        (heterogeneidad, costes de transacci\'on, negociaci\'on). Nivel ``natural''
        t\'ipico $\sim 5\%$ (Wheaton 1990).
  \item \textbf{Inactive / voluntary vacancy ($i$).}
        El propietario elige no participar. Dos mecanismos:
        \begin{itemize}
          \item \textbf{Opci\'on real / incertidumbre:} con precios inciertos,
                esperar tiene valor (Cunningham 2006).
          \item \textbf{Holding estrat\'egico} bajo regulaci\'on pro-inquilino:
                si no puedes ajustar renta ni terminar contratos f\'acilmente,
                puede ser racional sacar la unidad del alquiler
                (Gabriel \& Nothaft 2001).
        \end{itemize}
\end{enumerate}
Total medido: $V = v + i$.
\end{modelo}

\begin{resumenbox}
Grandes poseedores / inversores son los agentes para los que la
\emph{outside option} de dejar vac\'io ($b_j$) es sistem\'aticamente alta:
plusval\'ia, cartera diversificada, menor urgencia de cash-flow.
El tax apunta a ese margen.
\end{resumenbox}

\section{Modelo de participaci\'on (Seg\'u sobre Desgranges \& Wasmer)}

\begin{intuicion}
Imagina que cada due\~no compara: ``¿me compensa meter el piso en el mercado
o me quedo fuera?'' El impuesto hace m\'as caro quedarse fuera. Quien de verdad
oferta (y lo demuestra) puede estar exento --- no castig\'ais la fricci\'on honesta.
\end{intuicion}

\begin{modelo}
Setup: equilibrio parcial, corto plazo, stock fijo.
\begin{itemize}
  \item Condici\'on de participaci\'on: $rV_j \ge b_j - t$.
  \item Subir $t$ eleva el umbral $b^{*}$: menos unidades inactivas.
\end{itemize}
\textbf{Predicciones:}
\begin{itemize}
  \item $i \downarrow$ (cae el stock inactivo)
  \item $v$ puede subir un poco (entran m\'as al matching)
  \item $V = v+i \downarrow$ neto (domina la ca\'ida de $i$)
  \item Corto plazo: presi\'on a la baja en rentas
  \item Largo plazo: posible efecto adverso sobre nueva construcci\'on
        (limitaci\'on del modelo ST --- admitirlo)
\end{itemize}
\end{modelo}

\begin{frase}
El impuesto no castiga fricciones; castiga \textbf{no participar}.
\end{frase}

\section{Tres justificaciones de intervenci\'on}

\begin{center}
\renewcommand{\arraystretch}{1.25}
\begin{tabularx}{\textwidth}{@{}>{\bfseries\raggedright\arraybackslash}p{2.6cm}
  >{\raggedright\arraybackslash}X
  >{\raggedright\arraybackslash}X@{}}
\toprule
Canal & L\'ogica & C\'omo lo usamos \\
\midrule
Equity / asequibilidad &
Con demanda alta y oferta inel\'astica, mucha vacancy es parad\'ojica &
``No hace falta esperar a que se construya: hay viviendas que ya existen'' \\
Eficiencia / externalidades &
Concentraciones de vac\'ios $\rightarrow$ $\downarrow$ valor de vecinos, inseguridad, cohesi\'on &
Caso Pigouviano: el vac\'io impone coste social no internalizado \\
Second-best / place-based &
Distorsiones previas (protecci\'on inquilino, real options) justifican un tax local &
En ES: Ley de vivienda + plusval\'ias $\rightarrow$ inactive stock plausible \\
\bottomrule
\end{tabularx}
\end{center}

\begin{dato}
Timing argument (frente a construction-led), cita Seg\'u (2020):\\[0.3em]
\emph{``Focusing on vacancy reduction may be a faster and easier way to increase
residential capacity than promoting the construction of new housing units.''}
\end{dato}

% ══════════════════════════════════════════════════════════════
\chapter{Evidencia estrella: Seg\'u (2020) --- Francia}
\label{ch:segu}

\section{Qu\'e hizo Francia (TLV 1999) --- plantilla de dise\~no}

\begin{intuicion}
Francia grav\'o pisos vac\'ios \textbf{mucho tiempo} (años), en ciudades grandes
con desequilibrio, con tipo creciente, y eximi\'o a quien demostraba que
intentaba alquilar/vender o reformar. Eso es el blueprint.
\end{intuicion}

\begin{center}
\renewcommand{\arraystretch}{1.2}
\begin{tabularx}{\textwidth}{@{}llX@{}}
\toprule
\textbf{Par\'ametro} & \textbf{TLV Francia} & \textbf{Implicaci\'on para nosotros} \\
\midrule
Definici\'on & Vac\'io $\ge$ \textbf{2 a\~nos} & Ataca long-term / especulaci\'on, no rotaci\'on \\
Tipo & 10\% $\rightarrow$ 12.5\% $\rightarrow$ 15\% valor locativo & M\'as caro cuanto m\'as se alarga \\
\'Ambito & Urbano $>200$k con desequilibrio & Focalizar mercados tensionados \\
Exenci\'on & Oferta alquiler/venta; obras mayores & No castigar matching / reforma \\
Resultado & Vacancy $\downarrow$ \textbf{0.8\,pp $\approx$ 13\%} & Ancla emp\'irica \#1 \\
Heterogeneidad & Efecto $\sim$50\% mayor donde hab\'ia m\'as vac\'ios & Targeting de hotspots \\
Margen & Concentrado en \textbf{long-term} & Encaja anti-especulaci\'on \\
Destino & Mayor\'ia $\rightarrow$ \textbf{primary residences} & No solo relabel a 2\textsuperscript{a} \\
Construcci\'on & Sin efecto significativo a medio plazo & Contra el miedo ``mata la oferta'' \\
\bottomrule
\end{tabularx}
\end{center}

\section{Claims seguros vs cuidadosos}

\begin{dato}[title={Claims seguros (decir en presentaci\'on)}]
\begin{itemize}
  \item Hay evidencia cuasi-experimental de que un vacancy tax reduce vacancy.
  \item El efecto es mayor en long-term vacancy y en municipios con m\'as vac\'ios previos.
  \item Las unidades tienden a pasar a residencia principal.
\end{itemize}
\end{dato}

\begin{cuidado}
Seg\'u \textbf{no} estima un efecto grande y limpio de $\downarrow$ rentas a corto
plazo en todos los specs. Mejor claim: ``libera stock residencial / reduce
inactive housing'', no ``bajará el alquiler un $X\%$ mañana''.
Extrapolaci\'on Francia $\rightarrow$ Espa\~na no es autom\'atica (medici\'on, enforcement, 2\textsuperscript{a}s).
\end{cuidado}

\section{Catalunya / Espa\~na ya tienen piezas del dise\~no}

\begin{intuicion}
No inventamos la rueda. Catalunya (desde 2015) grava vac\'ios $\ge 2$ a\~nos a
personas jur\'idicas (y grandes tenedores). El IBI +50\% a desocupados existe
en teor\'ia estatal, pero hist\'oricamente el enforcement fall\'o por falta de
definici\'on operable de ``desocupado''.
\end{intuicion}

\begin{dato}
\begin{itemize}
  \item Espa\~na (Censo 2011, IEB): vacancy $\sim$\textbf{12.3\%} en ciudades $>50$k
        vs ``natural'' $\sim 5\%$.
  \item Impuesto CAT: recaudaci\'on \textbf{2.44\,M€} (2023) $\rightarrow$ \textbf{7.01\,M€} (2024).
  \item Nuestra posici\'on puede \textbf{generalizar y endurecer}: umbrales por
        n\'umero de viviendas / personas jur\'idicas / carteras.
\end{itemize}
\end{dato}

% ══════════════════════════════════════════════════════════════
\chapter{Airbnb, segundas residencias y desv\'io de stock}
\label{ch:airbnb}

\section{Garcia-L\'opez et al.\ (2020) --- Barcelona}

\begin{intuicion}
Si un piso pasa de alquiler largo a tur\'istico permanente, no ``faltan
ladrillos'': se ha \textbf{desviado} stock residencial. Eso sube rentas y precios.
Vacancy tax y reglas STR son la misma familia de instrumentos.
\end{intuicion}

\begin{dato}
Evidencia causal BCN:
\begin{itemize}
  \item Barrio medio: rentas \textbf{+1.9\%}, precios transacci\'on \textbf{+4.6\%}.
  \item Top decil Airbnb: rentas \textbf{+7\%}, precios hasta \textbf{+17\%}.
  \item En 2015, $\sim$2\% del stock en Airbnb; $\sim$7\% de las unidades en alquiler.
  \item Con $\sim$10 d\'ias/mes de Airbnb se igualaba la renta mensual de largo plazo
        --- incentivo potente a desviar.
\end{itemize}
\end{dato}

\textbf{C\'omo lo usamos sin abandonar vacancy tax:}
\begin{enumerate}
  \item El problema no es solo oferta nueva: es \textbf{asignaci\'on del stock}.
  \item Vacancy tax + fiscalidad STR son complementarios.
  \item Grandes poseedores / hosts multi-listing = margen intensivo
        (misma narrativa anti-concentraci\'on).
\end{enumerate}

\section{BdE Box 4: holgura contable del stock}

\begin{intuicion}
El Banco de Espa\~na dice, en esencia: s\'i hay gap de construcci\'on,
\textbf{pero} si moviliz\'ais 2\textsuperscript{a}s + tur\'istico/seasonal + vac\'ias,
Madrid y Barcelona podr\'ian absorber cientos de miles de hogares.
\end{intuicion}

\begin{dato}
\begin{itemize}
  \item Gap acumulado 2022--24 (completions vs hogares): \textbf{400--450 mil} viviendas.
  \item Capacidad si se moviliza stock no residencial:
        $\sim$\textbf{390 mil} hogares (Madrid), $\sim$\textbf{310 mil} (Barcelona).
  \item $+\,\sim 50\,000$ viviendas tur\'isticas en 2024 vs 2023.
  \item Stock de extranjeros no residentes: $>$\textbf{500 mil} viviendas.
  \item Compraventas de no residentes: $\sim$\textbf{8.4\%} del total (2024).
\end{itemize}
\end{dato}

\begin{resumenbox}
Equipo construction-led: ``hay que construir 400k''.\\
Nosotros: construir es lento; parte del shortage es \textbf{uso del stock};
vacancy tax ataca ese margen, m\'as barato y m\'as r\'apido.
\end{resumenbox}

\section{Ryan-Collins / UCL: financiarizaci\'on}

\begin{intuicion}
Si tratamos la vivienda como activo financiero (homes as assets), el vac\'io
especulativo es un s\'intoma. Subsidiar compradores (estilo Help-to-Buy) con
oferta r\'igida \textbf{infla} precios. Mejor: gravar demanda especulativa / 2\textsuperscript{a}s
+ vivienda p\'ublica asequible.
\end{intuicion}

Esto encaja exactamente con: vacancy tax + uso de recaudaci\'on + complemento
STR/2\textsuperscript{a}s.

% ══════════════════════════════════════════════════════════════
\chapter{Datos oficiales de Espa\~na (diagn\'ostico)}
\label{ch:datos}

Esta parte no es ``paper'': es lo que pod\'eis citar como \textbf{institucional}
(BdE, INE, MIVAU, Catastro, BOE).

\section{Censo 2021 --- intensidad de uso (INE)}

\begin{cuidado}
La definici\'on 2021 es el\'ectrica (sin contrato o consumo $<$ umbral $\approx$ 15 d\'ias/año).
\textbf{No} es la taxonom\'ia principal/secundaria/vac\'ia del Censo 2011.
No dig\'ais ``subi\'o de 13.7\% a 14.4\%'' sin explicar el cambio de m\'etodo.
\end{cuidado}

\begin{center}
\renewcommand{\arraystretch}{1.15}
\begin{tabular}{@{}lrrrr@{}}
\toprule
\textbf{\'Ambito} & \textbf{Totales} & \textbf{Vac\'ias} & \textbf{\% vac\'ia} & \textbf{Uso espor\'adico} \\
\midrule
España & 26\,623\,708 & 3\,828\,307 & \textbf{14.4\%} & 9.5\% \\
Madrid capital & 1\,533\,221 & 96\,921 & 6.3\% & --- \\
Barcelona ciudad & 808\,751 & 75\,488 & 9.3\% & --- \\
Alicante prov. & 1\,350\,006 & 209\,487 & 15.5\% & --- \\
Málaga prov. & 996\,557 & 152\,843 & 15.3\% & --- \\
\bottomrule
\end{tabular}
\end{center}

\begin{cuidado}
En ciudades $>250$k hay relativamente \textbf{poca} fracci\'on del total de vac\'ias
(BdE: solo $\sim$400k de $\sim$4M). El claim anti-especulaci\'on en centros tensionados
debe apoyarse tambi\'en en \textbf{2\textsuperscript{a}s + STR + no residentes + multi-propietarios},
no solo en ``hay millones de vac\'ias en Madrid centro''.
\end{cuidado}

\section{Censo 2011 (referencia IEB / Tempus3)}

\'Util para comparar con la literatura (IEB $\sim$12.3\% ciudades grandes) y con
Francia hist\'orica. Espa\~na 2011: tasa vac\'ia \textbf{13.65\%}
(3.44M de 25.2M viviendas).

\section{Construcci\'on vs hogares (timing)}

\begin{dato}
\begin{itemize}
  \item Terminadas 2024 (MIVAU): $\sim$\textbf{101 mil}.
  \item Creaci\'on neta de hogares: $\sim$\textbf{180 mil} (2024); promedio
        $\sim$\textbf{245 mil}/año en 2022--23 (BdE / INE ECP).
  \item Conclusi\'on de debate: el flujo nuevo no cierra el gap pronto;
        el stock ya existente s\'i puede moverse con incentivos.
\end{itemize}
\end{dato}

\section{Grandes tenedores (Catastro 2025)}

\begin{intuicion}
``Anti-grandes poseedores'' no es eslogan: el Catastro cuenta cu\'antos titulares
tienen muchas viviendas. Ah\'i est\'a la base para umbrales de dise\~no.
\end{intuicion}

\begin{center}
\renewcommand{\arraystretch}{1.15}
\begin{tabular}{@{}lrr@{}}
\toprule
\textbf{Tramo viviendas/titular} & \textbf{Titulares} & \textbf{Viviendas} \\
\midrule
1 & 12\,941\,416 & 12\,941\,416 \\
2--5 & 3\,974\,244 & 9\,515\,931 \\
6--10 & 104\,314 & 737\,462 \\
11--25 (``gran tenedor'' Ley 12/2023) & 21\,515 & 316\,766 \\
26--100 & 4\,539 & 195\,321 \\
$>100$ & 843 & 528\,868 \\
\midrule
\textbf{$\ge 11$ (suma ``grandes'')} & \textbf{26\,897} & \textbf{1\,040\,955 ($\sim$4.3\%)} \\
\bottomrule
\end{tabular}
\end{center}

\textbf{Marco legal \'util:}
\begin{itemize}
  \item Ley 12/2023: gran tenedor $>10$ inmuebles residenciales (o $>1500\,\mathrm{m}^{2}$);
        en tensionadas CCAA puede bajar a $\ge 5$.
  \item Recargo IBI desocupadas: $>2$ a\~nos, m\'inimo 4 viviendas/titular; hasta $+150\%$.
  \item Impuesto CAT vac\'ias: personas jur\'idicas + grandes tenedores ($>15$);
        recaudaci\'on finalista a vivienda protegida.
\end{itemize}

\section{Precios y alquileres (INE, contexto)}

\begin{itemize}
  \item IPV (compraventa): variaci\'on anual reciente $\sim$\textbf{+12\%} nacional;
        Madrid $\sim$\textbf{+13\%}.
  \item IPVA (alquiler): \'indice 2015$=100$ $\rightarrow$ 2024$=122.9$ ($\sim$\textbf{+23\%} acumulado).
  \item VTE (tur\'isticas): stock nacional reciente del orden de \textbf{341k--403k}
        seg\'un corte; BdE destaca el salto 2023$\rightarrow$2024.
\end{itemize}

\section{Extrapolabilidad Espa\~na--Francia}

\begin{dato}
\begin{itemize}
  \item INSEE: $\sim$\textbf{7.6\%} logements vacants (m\'etropole 2025).
  \item Eurostat Census 2021 \emph{unoccupied}: ES \textbf{30.4\%} vs FR \textbf{17.9\%}.
  \item Claim seguro: ES tiene m\'as stock no-usual-residence que FR.
  \item Claim cuidadoso: definiciones \textbf{no son id\'enticas} (el\'ectrica ES vs
        recensement/fiscal FR). No mezclar 14.4\% con 7.6\% como si fueran el mismo objeto.
\end{itemize}
\end{dato}

\section{Claims listos (solo institucionales)}

\begin{enumerate}
  \item ``El BdE cifra el d\'eficit acumulado 2022--24 entre terminaciones y hogares
        en 400--450 mil viviendas.''
  \item ``Si se movilizasen 2\textsuperscript{a}s, usos tur\'isticos/seasonal y vac\'ias,
        Madrid $\sim$390 mil hogares y Barcelona $\sim$310 mil (BdE 2024).''
  \item ``Censo 2021 (consumo el\'ectrico): 3.83 millones de vac\'ias (14.4\%);
        Madrid capital $\sim$6.3\%, Barcelona ciudad $\sim$9.3\%.''
  \item ``M\'as de 500 mil viviendas de extranjeros no residentes; compras de no residentes
        $\sim$8.4\% (BdE).''
  \item ``Catastro: $\sim$27 mil titulares con m\'as de 10 viviendas residenciales,
        $\sim$1.04 millones de inmuebles ($\sim$4\% del stock).''
  \item ``Catalunya ya grava vac\'ios $\ge 2$ a\~nos a PJ/grandes tenedores;
        recaud\'o 7\,M€ en 2024 (Idescat).''
  \item ``En 2024 se terminaron $\sim$101 mil viviendas (MIVAU) frente a una creaci\'on
        de hogares mucho mayor --- la v\'ia stock es el complemento de timing.''
\end{enumerate}

\begin{cuidado}[title={Lo que NO hay como dato oficial}]
Elasticidades, da\~no marginal Pigou, efecto causal del tax en ES, y vacancy
``natural'' $\sim$5\% (literatura). Para eso se citan papers del source pack
como evidencia, o se dejan como supuestos de dise\~no.
\end{cuidado}

% ══════════════════════════════════════════════════════════════
\chapter{Dise\~no de nuestra propuesta}
\label{ch:diseno}

\section{N\'ucleo de la medida}

\begin{frase}
Impuesto sobre vacancy \textbf{prolongada}, con umbrales/\'enfasis en
multi-propietarios y grandes poseedores; recaudaci\'on hacia vivienda
asequible / limited-profit; \textbf{complementaria} a (no sustituta de) oferta nueva.
\end{frase}

\section{Lecciones de welfare (IEB + Seg\'u)}

\begin{center}
\renewcommand{\arraystretch}{1.25}
\begin{tabularx}{\textwidth}{@{}>{\raggedright\arraybackslash}p{4cm} X@{}}
\toprule
\textbf{Factor} & \textbf{Efecto / implicaci\'on} \\
\midrule
Qui\'en ocupa las ex-vac\'ias &
Si pasan a residencia primaria $\rightarrow$ gana asequibilidad; si solo se relabelan como 2\textsuperscript{a} $\rightarrow$ no \\
Uso de la recaudaci\'on &
Reinvertir en vivienda social/asequible amplifica el efecto (ej.\ Catalunya) \\
Distorsi\'on LT a inversi\'on &
Teor\'ia: puede $\downarrow$ construction; empiria francesa: no se ve a medio plazo \\
Calidad / obras &
Exenciones por reforma evitan castigar rehabilitaci\'on \\
Medici\'on &
Sin definici\'on legal operable de ``vac\'io'', el tax muere (problema ES cl\'asico) \\
\bottomrule
\end{tabularx}
\end{center}

\section{Plantilla de dise\~no (para slides)}

Dedicad \textbf{al menos una slide} a esto --- convierte el slogan en
``government proposal'':

\begin{enumerate}[leftmargin=1.8em]
  \item \textbf{Definici\'on:} vac\'ia si sin uso residencial efectivo $\ge$ 2 a\~nos
        (o 1 a\~no en zonas tensionadas, alineable con TLV francesa actual);
        proxy: catastro + consumo el\'ectrico / agua + cruces empadronamiento.
  \item \textbf{Sujetos:} tipos m\'as altos / umbrales para personas jur\'idicas y
        titulares con $\ge N$ viviendas (inspiraci\'on CAT + Ley 12/2023).
  \item \textbf{Tipo:} creciente con la duraci\'on de la vacancy (l\'ogica francesa).
  \item \textbf{Exenciones:} oferta seria de alquiler/venta; obras mayores;
        herencia reciente / circunstancias personales acotadas.
  \item \textbf{\'Ambito:} mercados tensionados primero (Madrid, BCN, costas, islas).
  \item \textbf{Recaudaci\'on:} finalista a vivienda asequible / limited-profit.
  \item \textbf{Complementos} (misma familia, sin diluir el mensaje):
        recargo 2\textsuperscript{a}s / no residentes / STR permanente.
\end{enumerate}

\section{Extrapolaci\'on a ``anti-grandes poseedores'' (sin traicionar papers)}

\begin{enumerate}
  \item \textbf{Incidencia esperada:} el coste es m\'as alto para quien mantiene
        muchas unidades inactivas.
  \item \textbf{Heterogeneidad (Seg\'u):} m\'as impacto donde vacancy inicial es alta.
  \item \textbf{Personas jur\'idicas / umbrales (Catalunya):} no pegar al propietario
        de una sola vivienda en transici\'on.
  \item \textbf{Complemento STR/2\textsuperscript{a}s:} grandes operadores sobrerrepresentados
        en multi-listing y carteras de inversi\'on.
\end{enumerate}

\begin{cuidado}
Seg\'u \textbf{no} identifica causalmente ``fondos buitre'' vs peque\~nos.
Argumentadlo como \textbf{dise\~no de incidencia}, no como hallazgo del paper.
\end{cuidado}

\begin{cuidado}[title={Recaudaci\'on CAT $\approx$ 7\,M€}]
Es peque\~na. Argumentad \textbf{comportamiento + finalismo}, no
``financia sola el parque p\'ublico''.
\end{cuidado}

% ══════════════════════════════════════════════════════════════
\chapter{El otro lado: c\'omo no suicidarnos (y c\'omo girarlo)}
\label{ch:otro-lado}

\section{Greenaway-McGrevy \& Phillips (2023) --- Auckland upzoning}

\begin{intuicion}
Upzoning \textbf{s\'i} aumenta permisos/construcci\'on. No lo neg\'eis.
Negad que sea la \textbf{\'unica} soluci\'on y recordad el timing.
\end{intuicion}

\textbf{Respuesta parlamentaria:}
\begin{enumerate}
  \item Distinguir \textbf{permitted} vs \textbf{completed} homes.
  \item Timing: construcci\'on tarda a\~nos; vacancy tax act\'ua sobre stock ya existente.
  \item Complementariedad: ``No estamos en contra de construir; priorizamos movilizar stock
        \textbf{y} no tratamos upzoning como soluci\'on \'unica''.
  \item Auckland $\neq$ Espa\~na (turismo, 2\textsuperscript{a}s, financiarizaci\'on, elasticidad).
\end{enumerate}

\section{Asquith, Mast \& Reed --- pisos nuevos y rentas locales}

\begin{intuicion}
Edificios nuevos market-rate bajan rentas cercanas (5--7\%) absorbiendo hogares
de renta alta. Eso \textbf{confirma} que subir oferta efectiva baja rentas ---
exactamente la microeconom\'ia de liberar vac\'ios. Distinto instrumento, distinto horizonte.
\end{intuicion}

\begin{frase}
No conceder: ``entonces solo hay que construir''.\\
Conceder: ``la oferta importa; la m\'as r\'apida es la que ya est\'a vac\'ia''.
\end{frase}

\section{Jofre--Mart\'inez--Seg\'u (2023) --- rent control Catalunya}

\begin{itemize}
  \item Rentas $-4\%$ a $-6\%$; sin evidencia clara de $\downarrow$ supply de contratos en su ventana.
  \item \textbf{No es nuestro instrumento}, pero:
        (a) muestra que pol\'iticas ``anti-mercado puro'' pueden funcionar en ES bajo dise\~no fuerte;
        (b) rent control fuerte puede \textbf{aumentar} voluntary vacancy --- el tax es el
        complemento second-best.
  \item Si el parlamento propone solo rent control: ``sin movilizar vac\'ios, los caps pueden
        coexistir con stock inactivo''.
\end{itemize}

\section{Chay--Greenstone / Linden--Rockoff (lectures)}

Capitalizaci\'on: precios revelan WTP por amenities / riesgos.
Uso lateral: externalidades de vacancy (blight, inseguridad) se capitalizan en vecinos
$\rightarrow$ argumento Pigouviano. No forzar: no son papers de vacancy tax.

% ══════════════════════════════════════════════════════════════
\chapter{Mapa r\'apido: qu\'e paper para qu\'e claim}
\label{ch:mapa}

\begin{center}
\renewcommand{\arraystretch}{1.2}
\begin{longtable}{@{}p{5.2cm}p{4.2cm}p{3.5cm}@{}}
\toprule
\textbf{Claim} & \textbf{Fuente ancla} & \textbf{Tipo} \\
\midrule
\endhead
Vacancy voluntaria es racional, no ``error'' & Seg\'u 2020; IEB 2/2018 & Modelo + teor\'ia \\
Vacancy por encima de la ``natural'' justifica intervenci\'on & IEB ($\sim$12\% vs $\sim$5\%) & Benchmark \\
Vacancy tax \textbf{reduce} vacancy (esp.\ long-term) & Seg\'u: $-13\%$ & Causal DiD+PSM \\
Unidades vuelven a uso residencial primario & Seg\'u 2020 & Mecanismo \\
M\'as r\'apido movilizar stock que construir & Seg\'u; slides BdE & Timing \\
Externalidades de vac\'ios $\rightarrow$ eficiencia & Seg\'u intro (Lee, Glaeser\ldots) & Teor\'ia \\
Distorsiones previas justifican tax place-based & Kline \& Moretti via Seg\'u & Second-best \\
Airbnb/STR desv\'ia stock $\rightarrow$ rentas/precios $\uparrow$ & Garcia-L\'opez et al.\ 2020 & Causal ES \\
Demanda no residencial aprieta mercados ES & BdE Box 4 & Datos oficiales \\
Movilizar 2\textsuperscript{a}s+tur\'istico+vac\'ias libera capacidad & BdE Box 4 & Contabilidad \\
Impuestos a demanda especulativa + vivienda p\'ublica & Ryan-Collins / UCL & Pol\'itica \\
Subsidios a demanda pueden inflar precios & Seg\'u lit.\ review & Contra-alternativa \\
Construcci\'on nueva baja rentas locales & Asquith et al. & Complementariedad \\
Upzoning aumenta permisos & Greenaway-McGrevy \& Phillips & Distinguir permisos vs uso \\
\bottomrule
\end{longtable}
\end{center}

% ══════════════════════════════════════════════════════════════
\chapter{Frases listas y trampas}
\label{ch:frases}

\section{Seis frases con ancla}

\begin{enumerate}[leftmargin=1.6em, itemsep=0.6em]
  \item ``En Francia, la primera evaluaci\'on cuasi-experimental de un vacancy tax
        encuentra una ca\'ida del \textbf{13\%} en la tasa de vac\'ios, concentrada en
        vacancy de largo plazo, con desplazamiento hacia residencias principales.''
        --- Seg\'u \emph{JPubE} 2020.
  \item ``Con oferta casi fija a corto plazo, gravar el vac\'io aumenta el coste de
        oportunidad de no alquilar y, en el modelo de search-and-matching, reduce
        el stock inactivo.'' --- Seg\'u / Desgranges--Wasmer.
  \item ``España lleva a\~nos con vacancy urbana por encima de la tasa `natural'
        ($\sim$5\% vs $\sim$12\% en ciudades grandes, Censo 2011).'' --- IEB 2/2018.
  \item ``Airbnb en Barcelona no es anecd\'otico: en barrios del top decil, rentas +7\%
        y precios hasta +17\%.'' --- Garcia-L\'opez et al.\ \emph{JUE} 2020.
  \item ``El BdE cuantifica que movilizar 2\textsuperscript{a}s, usos tur\'isticos/seasonal
        y vac\'ios liberar\'ia cientos de miles de hogares potenciales en Madrid y Barcelona.''
        --- Box 4, 2024.
  \item ``Los subsidios a la demanda en mercados r\'igidos tienden a capitalizarse en precios;
        mejor gravar demanda especulativa y reinvertir en vivienda asequible.''
        --- lit.\ allowances en Seg\'u + Ryan-Collins/UCL.
\end{enumerate}

\section{Lo que \emph{no} debemos decir}

\begin{cuidado}[title={Trampas}]
\begin{itemize}
  \item ``Seg\'u demuestra que bajan los alquileres un $X\%$ en Espa\~na.''
        (Paper Francia; foco principal = vacancy rate.)
  \item ``Upzoning no funciona.''
        (Auckland dice que s\'i aumenta construcci\'on; atacar exclusividad, no existencia.)
  \item ``Asquith prueba que construir no ayuda.''
        (Prueba lo contrario a escala local.)
  \item ``Todos los vac\'ios son especulaci\'on de fondos.''
        (Parte es fricci\'on; por eso umbral de a\~nos + exenciones.)
  \item ``El tax no tiene costes.''
        (Teor\'ia LT: posible $\downarrow$ inversi\'on; admitir trade-off + evidencia media plazo FR + reinversi\'on.)
  \item Mezclar Censo 2011 y 2021 como si la tasa fuera comparable sin matices.
  \item Usar recaudaci\'on catalana de 7\,M€ como si financiara sola el parque p\'ublico.
\end{itemize}
\end{cuidado}

% ══════════════════════════════════════════════════════════════
\chapter{Problem sets del curso --- checklist operativo}
\label{ch:ps}

\begin{center}
\renewcommand{\arraystretch}{1.2}
\begin{tabularx}{\textwidth}{@{}c>{\raggedright\arraybackslash}p{4cm} X@{}}
\toprule
\textbf{PS} & \textbf{Tema} & \textbf{Para nuestra postura} \\
\midrule
PS1 & Elasticidades, CV/EV & Bienestar de inquilinos ante cambio de rentas/precios \\
PS2 & Edgeworth, SFTWE, utilitarista vs Rawls & Equity slide; ``Pareto no basta'' \\
PS3 & Emp\'iricos & Disciplina causal al citar Seg\'u / BdE \\
PS4 & Externalidades + Pigou + Weitzman & \textbf{Plantilla formal} del argumento de eficiencia \\
PS5 & Public goods / free riding & Solo si habl\'ais de provisi\'on p\'ublica con recaudaci\'on \\
PS6 & Tax-transfer + median voter & Econom\'ia pol\'itica del apoyo a gravar poseedores \\
PS7 & CBA housing + discounting + distribution & Timing stock vs new build; impactos sociales \\
PS8 & Repaso / MCQ & Pigou, Coase, Samuelson --- vocabulario oral \\
\bottomrule
\end{tabularx}
\end{center}

% ══════════════════════════════════════════════════════════════
\chapter{Cierre: c\'omo montar los 10 minutos}
\label{ch:cierre}

\begin{intuicion}[title={Estructura sugerida de la presentaci\'on}]
\begin{enumerate}
  \item \textbf{Problema (1'):} crisis ES = demanda + stock mal usado (BdE 400--450k
        \emph{y} 310k/390k movilizables).
  \item \textbf{Fallo (1.5'):} inactive vacancy + externalidades + equity
        (curso A1--A4 en lenguaje simple).
  \item \textbf{Modelo (1.5'):} $V=v+i$; el tax baja $i$ (participaci\'on $b_j-t$).
  \item \textbf{Evidencia (2'):} Seg\'u $-13\%$; Airbnb BCN; Catastro grandes tenedores.
  \item \textbf{Propuesta (2.5'):} definici\'on, tipos, umbrales, exenciones,
        enforcement, uso de ingresos.
  \item \textbf{Vs construction (1'):} complementarios; nosotros priorizamos timing del stock.
  \item \textbf{Cierre (0.5'):} winners/losers + voto.
\end{enumerate}
\end{intuicion}

\begin{resumenbox}
Si solo record\'ais tres cosas:\\[0.4em]
(1) Hay vac\'ios voluntarios racionales --- un precio los reduce.\\
(2) Francia lo midió: $-13\%$, hacia residencia principal.\\
(3) En Espa\~na el BdE ya cuantifica la holgura del stock; nosotros ponemos el incentivo.
\end{resumenbox}

\vspace{1.2cm}
{\color{azulmedio}\rule{\textwidth}{1pt}}

\vspace{0.6cm}
\noindent\small\color{azuloscuro}
\textbf{Fuentes principales del dossier:}\\[0.3em]
Seg\'u (2020) \emph{JPubE}; IEB Report 2/2018; Garcia-L\'opez et al.\ (2020) \emph{JUE};
BdE Informe Anual 2024 Box~4; Ryan-Collins/UCL 2024; Greenaway-McGrevy \& Phillips (2023);
Asquith--Mast--Reed; Jofre--Mart\'inez--Seg\'u (2023); INE Censo 2021 (tpx~59531);
MIVAU; DG Catastro; BOE Ley 12/2023 y Ley 14/2015; material de curso Weeks~1--12 + PS1--8.\\[0.5em]
Archivos locales: \texttt{docs/arguments-from-papers.md}, \texttt{docs/course-arguments.md},
\texttt{docs/ine-data-from-mcp.md}, \texttt{research/briefs/research-brief-datos-oficiales-faltantes.md},
\texttt{references/}.

\end{document}
